\documentclass[11pt]{amsart}

\usepackage[T1]{fontenc}
\usepackage{lmodern}
\usepackage{microtype}
\usepackage{mathtools,amssymb,amsthm}
\usepackage[hidelinks]{hyperref}

\hypersetup{
  pdftitle={Unimodal unicyclic graphs with several breaks of log-concavity}
}

\newtheorem{theorem}{Theorem}
\newtheorem{lemma}[theorem]{Lemma}
\newtheorem{corollary}[theorem]{Corollary}
\theoremstyle{definition}
\newtheorem*{problemLK}{Problem (multi-break unicyclic families)}
\theoremstyle{remark}
\newtheorem*{remark}{Remark}

\setlength{\emergencystretch}{3em}

\newcommand{\Ip}{I}
\newcommand{\albr}{\alpha}
\newcommand{\TG}{T_{G}}

\title[Unicyclic graphs with several breaks]
{Unimodal unicyclic graphs with several breaks of log-concavity}

\date{}

\subjclass[2020]{05C31, 05C69, 05C05}
\keywords{independence polynomial, log-concavity, unimodality, unicyclic graph,
one-edge closure, pattern graph, counterexample}

\begin{document}

\begin{abstract}
For a graph $G$ write $\Ip(G;x)=\sum_{k\ge0}i_kx^k$, where $i_k$ is the number of
independent $k$-sets of $G$, and call an index $1\le k\le\albr(G)-1$ a \emph{break}
of $G$ if $i_k^2<i_{k-1}i_{k+1}$. Levit and Kadrawi
\cite[Problem 6.6]{LevitKadrawi} ask whether one-edge closure of a
multi-break or pattern-graph tree can produce a \emph{unicyclic} graph whose
independence polynomial is unimodal and has more than one break, or even several
\emph{consecutive} breaks. We exhibit unimodal unicyclic graphs of both kinds.
The tree $\TG{}_{,2,6}$ closed at one nonedge gives a unimodal unicyclic graph on
$82$ vertices with breaks exactly $\{40,42\}$, and a pattern graph on $58$ vertices
closed at one nonedge gives breaks exactly $\{28,29\}$, two \emph{consecutive} breaks.
One closure of $\TG{}_{,7,6}$, on $282$ vertices, is unimodal with \emph{seven} breaks
while its tree has six, so one-edge closure can also increase the number of breaks; in
the five other witnesses below, including the two just named, the closure preserves the
break set of the tree it closes. All arithmetic is exact integer arithmetic.
\end{abstract}

\maketitle

\section{The problem}

Let $G$ be a finite simple graph, $i_k=i_k(G)$ the number of independent sets of $G$
of size $k$, $\albr(G)=\max\{k:i_k>0\}$ the independence number, and
\[
  \Ip(G;x)=\sum_{k=0}^{\albr(G)}i_kx^k
\]
the independence polynomial. The sequence $(i_k)_{k=0}^{\albr}$ is \emph{unimodal} if
it is nondecreasing up to some index and nonincreasing thereafter, and
\emph{log-concave} if $i_k^2\ge i_{k-1}i_{k+1}$ for $1\le k\le\albr-1$. We call an
index $k$ in that range a \emph{break} of $G$ when the inequality fails,
\begin{equation}\label{eq:break}
  i_k^2<i_{k-1}i_{k+1},
\end{equation}
and write $B(G)\subseteq\{1,\dots,\albr(G)-1\}$ for the set of breaks. A graph is
\emph{unicyclic} if it is connected and $|E|=|V|$; a \emph{one-edge closure} of a tree
$T$ is $T+e$ for a nonedge $e$ of $T$, which is exactly a unicyclic graph.

The question addressed here is Problem~6.6 of Levit and Kadrawi
\cite{LevitKadrawi}.

\begin{problemLK}
Starting from Bautista--Ramos's multi-break families
\textup{\cite{BautistaRamos}} and from the newer pattern-graph families with two and
three consecutive breaks \textup{\cite{BRGG}}, decide whether one-edge closure can
produce unicyclic graphs whose independence polynomials fail log-concavity at more
than one index, or even at several consecutive indices, while remaining unimodal.
\end{problemLK}

Nothing below is a tree-side result: the trees we close come from the families and the
calculus of \cite{BautistaRamos,BRGG}, and the tree break sets we quote are recomputed
only as baselines. The graph $H_1$ below is one of the closures of $\TG{}_{,2,6}$
already computed in \cite{LevitKadrawi}; what we record for it, and for the other five
witnesses, is the number and the position of the breaks. The theorem below is a
statement about six explicit graphs and does not depend on how the problem above is
read.

\section{The graphs}

We use the pattern-graph calculus of Bautista--Ramos, Guill\'en-Galv\'an and
G\'omez-Salgado \cite{BRGG}. All graphs are rooted; the
root of $X$ is written $r(X)$.

\begin{itemize}
\item $P_t$ is the path on $t$ vertices, rooted at a leaf.
\item $Z_k(H)$ is a new vertex $v_0$ together with $k$ disjoint copies of $H$, with
      $v_0$ joined to the root of each copy; $r(Z_k(H))=v_0$.
\item $(G:H)_k$ is the disjoint union of $G$ and $Z_k(H)$ with the edge
      $r(G)\,v_0$ added; the root stays $r(G)$.
\item $(G:H)_k^{(m)}$ is the result of applying $(\,\cdot\,:H)_k$ to $G$ $m$ times in
      succession, the root staying $r(G)$ throughout. Hence
      \begin{equation}\label{eq:size}
        \bigl|(G:H)_k^{(m)}\bigr|=|G|+m\bigl(1+k|H|\bigr).
      \end{equation}
\item $S_{2,t}=(P_1:P_1)_1^{(t)}$: a centre with $t$ pendant paths of length two;
      $1+2t$ vertices.
\item $T_{1,l}=(P_1:P_2)_l^{(1)}$: a root $v$ joined to a vertex $v_0$ that carries
      $l$ pendant paths of length two; $2+2l$ vertices.
\item $\TG{}_{,m,t}=(P_2:S_{2,t})_3^{(m)}$, the Bautista--Ramos multi-break tree
      (written $GT_{m,t}$ in \cite{BautistaRamos,BRGG}): a root $v_0$ with one extra
      leaf and $m$ pendant blocks, each block a hub joined to the centres of three
      copies of $S_{2,t}$. By \eqref{eq:size}, $|\TG{}_{,m,t}|=m(4+6t)+2$.
\end{itemize}

\medskip
\noindent\textbf{Labelling.} Because our witnesses are named by the integer pair of
their closure edge, we fix the labelling that the construction induces: vertices are
numbered $0,1,2,\dots$ in the order in which the recipe
above creates them. Concretely, in $(G:H)_k^{(m)}$ the vertices of $G$ come first, in
$G$'s own order ($P_t$ being labelled $0,1,\dots,t-1$ along the path from its root);
then the $m$ blocks in order, each block contributing its hub first and then its $k$
copies of $H$ in order, each copy in $H$'s own order. In $S_{2,t}$, for instance, this
makes $0$ the centre and the $j$-th pendant path the pair $(2j-1,2j)$ for
$1\le j\le t$, with $2j-1$ the middle vertex and $2j$ the leaf. Every closure edge
below is also described in words, so no reader has to trust the numbering.

\medskip
\noindent Thus, explicitly:
\begin{itemize}
\item $H_1=\TG{}_{,2,6}+\{0,3\}$, joining the root $v_0=0$ to the centre $3$ of the
      first copy of $S_{2,6}$ in the first block, whose hub is $2$; the unique cycle
      is the triangle $0\,2\,3$.
\item $H^*=(P_1:S_{2,4})_2^{(3)}+\{0,3\}$, joining the root $v=0$ to the vertex $3$,
      a middle vertex of a pendant path of the centre $c=2$ of the first copy of
      $S_{2,4}$ in the first block, whose hub is $u=1$; the unique cycle is the
      $4$-cycle $0\,1\,2\,3$.
\item $H^{**}=(P_3:S_{2,4})_2^{(8)}+\{0,5\}$, joining the root $v=0$ of the base path
      $0\,1\,2$ to the middle vertex $5$ of a pendant path of the centre $4$, whose
      hub is $3$; the unique cycle is the $4$-cycle $0\,3\,4\,5$.
\item $H_A=(T_{1,3}:S_{2,4})_2^{(16)}+\{0,10\}$, joining the root $0$ of the base
      $T_{1,3}$ (labelled $0,\dots,7$) to the middle vertex $10$ of a pendant path of
      the centre $9$, whose hub is $8$; the unique cycle is the $4$-cycle
      $0\,8\,9\,10$.
\item $H_B=(T_{1,7}:S_{2,5})_2^{(13)}+\{1,72\}$, joining the vertex $1$ of the base
      $T_{1,7}$ (labelled $0,\dots,15$; $1$ is the neighbour of the root) to the
      middle vertex $72$ of a pendant path of the centre $63$ in the third block,
      whose hub is $62$; the unique cycle is the $5$-cycle $1\,0\,62\,63\,72$.
\item $H_7=\TG{}_{,7,6}+\{0,3\}$, the same closure edge as $H_1$ one block-count
      higher; the unique cycle is the triangle $0\,2\,3$.
\end{itemize}

\section{The witnesses}

\begin{theorem}\label{thm:main}
Each of the following six graphs is unicyclic, its independence polynomial is
unimodal with a unique mode, and its set of breaks is \emph{exactly} the set listed
--- not merely a subset of it.

\begin{center}
\renewcommand{\arraystretch}{1.15}
\begin{tabular}{llrrrl}
 & $H$ & $|V|=|E|$ & $\albr(H)$ & mode & $B(H)$\\
\hline
\textup{(a)} & $H_1$      & $82$  & $43$  & $i_{25}$ & $\{40,42\}$\\
\textup{(b)} & $H^*$      & $58$  & $31$  & $i_{18}$ & $\{28,29\}$\\
\textup{(c)} & $H^{**}$   & $155$ & $82$  & $i_{47}$ & $\{78,79,80\}$\\
\textup{(d)} & $H_A$      & $312$ & $164$ & $i_{95}$ & $\{162,163\}$\\
\textup{(e)} & $H_B$      & $315$ & $164$ & $i_{97}$ & $\{161,162,163\}$\\
\textup{(f)} & $H_7$      & $282$ & $148$ & $i_{88}$ & $\{135,137,139,141,143,145,147\}$
\end{tabular}
\end{center}
\end{theorem}

\begin{corollary}\label{cor:answer}
One-edge closure produces unimodal unicyclic graphs whose independence polynomials
fail log-concavity at more than one index --- at two indices in \textup{(a)}, and at
seven in \textup{(f)} --- and at several \emph{consecutive} indices --- two in
\textup{(b)} and three in \textup{(c)} and \textup{(e)}. Here \textup{(d)} and
\textup{(e)} are closures of $(T_{1,3}:S_{2,4})_2^{(16)}$ and
$(T_{1,7}:S_{2,5})_2^{(13)}$, members of the two pattern-graph families of
\textup{\cite{BRGG}}, and \textup{(a)} and \textup{(f)} closures of
Bautista--Ramos's $GT_{m,t}$ \textup{\cite{BautistaRamos}}.
\end{corollary}

\subsection*{The $58$-vertex witness, in full}
$H^*$ has $58$ vertices, so its whole independence polynomial fits on the page and
Theorem~\ref{thm:main}(b) can be checked with a calculator. Its coefficients
$i_0,\dots,i_{31}$ are
{\footnotesize
\begin{verbatim}
1, 58, 1595, 27709, 341818, 3190686, 23448526, 139359013, 682663978,
2794877442, 9661712907, 28414927226, 71476627753, 154326951878,
286560418561, 457792570649, 628535097664, 739676512678, 742815300664,
632430324122, 452340586723, 268386668712, 129808551984, 49924871616,
14717851602, 3137062048, 435476016, 31145353, 415590, 6304, 113, 1
\end{verbatim}}
\noindent
The sequence increases strictly up to $i_{18}=742815300664$ and decreases strictly
afterwards, so it is unimodal with a unique mode. Four multiplications settle the
break set at the top:
\begin{align*}
 k=27:&& i_{27}^2&=970033013494609 &\ge&& i_{26}i_{28}&=180979477489440 &&\text{(holds)},\\
 k=28:&& i_{28}^2&=172715048100 &<&& i_{27}i_{29}&=196340305312 &&\text{(break)},\\
 k=29:&& i_{29}^2&=39740416 &<&& i_{28}i_{30}&=46961670 &&\text{(break)},\\
 k=30:&& i_{30}^2&=12769 &\ge&& i_{29}i_{31}&=6304 &&\text{(holds)}.
\end{align*}
The remaining $26$ inequalities $1\le k\le26$ hold; they are large but immediate from
the list. So $B(H^*)=\{28,29\}$, two consecutive breaks.

In \texttt{graph6} format $H^*$ is the concatenation of the following five lines (no
spaces, $277$ characters, single-byte header since $58\le62$):
{\footnotesize
\begin{verbatim}
ylC_H?@G??g??@??_?_?@?@???G?G???E??????G???_??@???C????_
??A????@???C?????G??@??????C????@?????_?????_????O?????@
?????_??????K????????????C??????G??????G??????O??????@??
????A????????_?????@????????@??????C?????????G???????@??
??????O????????G???????A?????????C???????@??????????G
\end{verbatim}}

\subsection*{The decisive coefficients of the other five}
In each case the full coefficient list is determined by the construction of \S2; we
print the coefficients that decide the break set, and the neighbouring index at which
log-concavity holds, so that each entry of Theorem~\ref{thm:main} can be confirmed
without recomputing the whole polynomial.

\medskip
\noindent$H_1$, $\albr=43$: $i_{39}=78555205917$, $i_{40}=125984042$,
$i_{41}=622042$, $i_{42}=485$, $i_{43}=1$. Then
\begin{align*}
 i_{40}^2&=15871978838657764<i_{39}i_{41}=48864637399022514,\\
 i_{42}^2&=235225<i_{41}i_{43}=622042,
\end{align*}
while at the index between them log-concavity holds,
$i_{41}^2=386936249764\ge i_{40}i_{42}=61102260370$: the two breaks of $H_1$ are
\emph{not} consecutive.

\medskip
\noindent$H^{**}$, $\albr=82$: $i_{77}=1135175576719$, $i_{78}=3247342473$,
$i_{79}=9297742$, $i_{80}=51245$, $i_{81}=315$, $i_{82}=1$. Then
{\footnotesize
\begin{align*}
 i_{78}^2&=10545233136949755729, & i_{77}i_{79}&=10554569637034468498,\\
 i_{79}^2&=86448006298564,       & i_{78}i_{80}&=166410065028885,\\
 i_{80}^2&=2626050025,           & i_{79}i_{81}&=2928788730,
\end{align*}}
\noindent and in each of the three rows the left-hand value is the smaller, so $78$, $79$ and
$80$ are breaks, while $i_{81}^2=99225\ge i_{80}i_{82}=51245$ holds. The first
of the three is decided by a margin of $0.088\%$, which is why the arithmetic here is
exact integer arithmetic and never floating point.

\medskip
\noindent$H_A$, $\albr=164$: $i_{160}=14399685489601$, $i_{161}=15059609304$,
$i_{162}=12154773$, $i_{163}=9815$, $i_{164}=9$. Then
\begin{align*}
 i_{162}^2&=147738506681529<i_{161}i_{163}=147810065318760,\\
 i_{163}^2&=96334225<i_{162}i_{164}=109392957,
\end{align*}
while at the index below, log-concavity holds:
{\footnotesize
\begin{align*}
 i_{161}^2&=226791832389123364416, & i_{160}i_{162}&=175024908397494015573.
\end{align*}}

\medskip
\noindent$H_B$, $\albr=164$: $i_{159}=1836675801123505133$,
$i_{160}=946417052360718$, $i_{161}=367642616716$, $i_{162}=146769284$,
$i_{163}=137465$, $i_{164}=129$. Then
{\footnotesize
\begin{align*}
 i_{161}^2&=135161093625787682624656, & i_{160}i_{162}&=138904953140373090585912,\\
 i_{162}^2&=21541222725872656,        & i_{161}i_{163}&=50537992306864940,\\
 i_{163}^2&=18896626225,              & i_{162}i_{164}&=18933237636,
\end{align*}}
\noindent and again the left-hand value is smaller in each of the three rows, so $161$, $162$ and
$163$ are breaks, while $i_{160}^2\ge i_{159}i_{161}$ holds. Three consecutive breaks,
inside one of the two pattern-graph families of \cite{BRGG}.

\medskip
\noindent$H_7$, $\albr=148$. Its top fifteen coefficients $i_{134},\dots,i_{148}$ are
{\footnotesize
\begin{verbatim}
176034211309656718484117918890005878151
966322839077538885713606551341161815
5610322511779111055524711146550627
19499704262596590695297758206411
69973942869517833994033432381
174438847810507838032800383
444508352733368700464454
822213872739649589233
1559903682459718229
2122169611386407
3022643280657
2853892767
2959617
1535
1
\end{verbatim}}
\noindent
and the seven inequalities \eqref{eq:break} at
$k=135,137,139,141,143,145,147$ all fail while the six at
$k=136,138,140,142,144,146$ hold, giving alternating breaks
$B(H_7)=\{135,137,\dots,147\}$.

\section{Proof of Theorem~\ref{thm:main}}

Everything reduces to computing one independence polynomial per graph, and for a
unicyclic graph that is a two-line reduction to forests.

\begin{lemma}\label{lem:rec}
\textup{(i)} If $F$ is a forest with components $F_1,\dots,F_c$ then
$\Ip(F;x)=\prod_j\Ip(F_j;x)$, and for a tree rooted at $r$ the two polynomials
$A_v$ (independent sets of the subtree at $v$ containing $v$) and $B_v$ (those not
containing $v$) satisfy
\[
 A_v=x\prod_{w\ \mathrm{child\ of}\ v}B_w,\qquad
 B_v=\prod_{w\ \mathrm{child\ of}\ v}(A_w+B_w),\qquad
 \Ip=A_r+B_r .
\]
\textup{(ii)} If $H$ is unicyclic and $uv$ is an edge of its unique cycle, then
$H-uv$ is a tree and
\begin{equation}\label{eq:uni}
 \Ip(H;x)=\Ip(H-uv;x)-x^2\,\Ip\bigl((H-uv)-N[u]-N[v];x\bigr),
\end{equation}
both graphs on the right being forests.
\textup{(iii)} If $w$ lies on the unique cycle of $H$, then
\begin{equation}\label{eq:split}
 \Ip(H;x)=\Ip(H-w;x)+x\,\Ip(H-N[w];x),
\end{equation}
and again both graphs on the right are forests.
\end{lemma}

\begin{proof}
(i) is the standard rooted recursion: an independent set of the subtree at $v$ either
contains $v$, and then contains no child of $v$, or does not. (ii): independent sets
of $H$ are exactly the independent sets of $H-uv$ that do not contain both $u$ and
$v$; those containing both are the sets $\{u,v\}\cup S$ with $S$ independent in
$(H-uv)-N[u]-N[v]$, whence \eqref{eq:uni}. Deleting a cycle edge of a unicyclic graph
leaves a spanning tree, and deleting further vertices leaves a forest. (iii) is the
same dichotomy applied to $w$; deleting a vertex of the unique cycle destroys it, so
$H-w$ is a forest, and $H-N[w]$ is an induced subgraph of it.
\end{proof}

\begin{proof}[Proof of Theorem~\ref{thm:main}]
Each $H$ in the table is $T+e$ for the tree $T$ and nonedge $e$ named in \S2. A tree
on $n$ vertices has $n-1$ edges, so $|E(H)|=|V(H)|$, and adding an edge to a
connected graph keeps it connected: $H$ is unicyclic. The vertex counts follow from
\eqref{eq:size}: $|\TG{}_{,2,6}|=2(4+36)+2=82$ and $|\TG{}_{,7,6}|=7\cdot40+2=282$;
$|(P_1:S_{2,4})_2^{(3)}|=1+3(1+2\cdot9)=58$;
$|(P_3:S_{2,4})_2^{(8)}|=3+8\cdot19=155$;
$|(T_{1,3}:S_{2,4})_2^{(16)}|=8+16\cdot19=312$;
$|(T_{1,7}:S_{2,5})_2^{(13)}|=16+13\cdot23=315$. In each case the two ends of $e$ are
joined in $T$ by the path displayed in \S2, so the unique cycle is as stated.

For the polynomial, apply Lemma~\ref{lem:rec}: \eqref{eq:uni} with $uv$ the closure
edge itself expresses $\Ip(H;x)$ through two forest polynomials, each computed by the
rooted recursion (i) in exact integer arithmetic. Then $\albr(H)=\deg\Ip(H;x)$;
unimodality and the mode are read off the coefficient list; and each break is the
single integer comparison \eqref{eq:break}. For $H^*$ the resulting coefficient list
and the four decisive comparisons are printed in \S3 in full, so
Theorem~\ref{thm:main}(b) is verified on the page. For the remaining five graphs the
decisive comparisons are printed in \S3, and \emph{every} index $1\le k\le\albr-1$ was
checked; the accompanying program recomputes all six polynomials from the definitions
of \S2 by the two independent routes \eqref{eq:uni} and \eqref{eq:split} and confirms
that the break sets are exactly those listed.
\end{proof}

\section{Remarks}

\begin{remark}
In (a)--(e) of Theorem~\ref{thm:main} the closure preserves the break set of the
underlying tree exactly, with $\albr$ unchanged; in (f) it does not, since
$B(\TG{}_{,7,6})=\{135,137,139,141,143,145\}$, so there the closure adds the break at
the index $\albr(H_7)-1=147$. We do not claim that the number of breaks of a unimodal
unicyclic graph can be made arbitrarily large: seven, on $282$ vertices, is the largest
number exhibited here. No minimality is claimed for any of the six witnesses.
\end{remark}

\begin{remark}
The trees underlying $H^*$ and $H^{**}$ are pattern graphs in the sense of \cite{BRGG}
but are not rows of that paper's table of named families. Read strictly, so that ``the
newer pattern-graph families'' means those named families, the consecutive-break clause
is carried by $H_A$ and $H_B$, which do lie in them.
\end{remark}

\section{Reproduction}

\S2 determines all six graphs from their parameters alone, and Lemma~\ref{lem:rec}
determines the polynomials from the graphs. A program \texttt{verify.py}
accompanies this note, with a record of a run in
\texttt{verify.output.txt}; it needs only Python~3 and its standard library,
\begin{center}\texttt{python3 verify.py}\end{center}
and exits $0$ only if every check passes. It rebuilds each graph from the definitions
of \S2, computes each independence polynomial by both \eqref{eq:uni} and
\eqref{eq:split}, re-derives every integer printed above --- the orders, the
independence numbers, the modes, the coefficient lists, each printed product and the
break sets in full --- decodes the \texttt{graph6} string of \S3 and matches it against
the constructed $H^*$, and reproduces as controls published integers of
\cite{BautistaRamos,BRGG,KadrawiLevit}. It does not read the source e-print, so the
statement of the problem in \S1 lies outside its scope.

\begin{thebibliography}{9}

\bibitem{BautistaRamos}
C.~Bautista-Ramos,
\emph{Multiple breaks of log-concavity in the independence polynomials of trees},
arXiv:2511.00334, 2025.
\href{https://arxiv.org/abs/2511.00334}{arXiv:2511.00334}.

\bibitem{BRGG}
C.~Bautista-Ramos, C.~Guill\'en-Galv\'an, and P.~G\'omez-Salgado,
\emph{Linear recurrences for non-log-concave independence polynomials of trees},
Graphs and Combinatorics (2026);
\href{https://doi.org/10.1007/s00373-026-03054-4}{doi:10.1007/s00373-026-03054-4};
\href{https://arxiv.org/abs/2603.14204}{arXiv:2603.14204}.

\bibitem{KadrawiLevit}
O.~Kadrawi and V.~E. Levit,
\emph{The independence polynomial of trees is not always log-concave starting from
order $26$},
Ars Mathematica Contemporanea \textbf{25} (2025), no.~4, Paper P4.03;
\href{https://arxiv.org/abs/2305.01784}{arXiv:2305.01784}.

\bibitem{LevitKadrawi}
V.~E. Levit and O.~Kadrawi,
\emph{Closing trees into unicyclic counterexamples},
arXiv:2603.17114v1, 17 March 2026, 30 pp.
\href{https://arxiv.org/abs/2603.17114v1}{arXiv:2603.17114v1}.

\end{thebibliography}

\end{document}
